\documentclass[11pt]{article}

\usepackage{amsmath}
\usepackage{amssymb}
\usepackage{bm}
\usepackage{xcolor}
\usepackage{algorithm}
\usepackage{algpseudocode}
\usepackage{graphicx}
\usepackage{booktabs}
\usepackage[margin=1in]{geometry}

\definecolor{changeblue}{HTML}{1F5FBF}
\newcommand{\changed}[1]{\textcolor{changeblue}{#1}}
\newcommand{\DensityFB}{Density-FB}
\newcommand{\sg}{\operatorname{sg}}
\newcommand{\normalize}{\operatorname{normalize}}

\setlength{\parindent}{0pt}
\setlength{\parskip}{\baselineskip}

\title{Density-FB Derivation}
\date{}

\begin{document}

\maketitle


% ================================================================
\section{Preliminaries: Forward-Backward Representations}
\label{sec:preliminaries}
% ================================================================

\paragraph{Notation.}
The offline dataset induces a distribution $\mathcal D$ over transitions $(s,a,s')$, with state-action marginal $\mathcal D_{SA}$ and next-state marginal $\rho$. We assume $0<\gamma<1$ and use $\rho$ as the reference distribution over future states. Let $p_{\mathcal T}$ denote a training-task distribution over pairs $(z,r_z)$, where $z$ indexes a policy and $r_z:\mathcal S\rightarrow\mathbb R$ is its associated task reward. We write $p_Z$ for the $z$-marginal of $p_{\mathcal T}$ and $\pi_z(a\mid s)=\pi(a\mid s,z)$.

For a policy $\pi$, define the discounted successor measure
\begin{equation}
    M^\pi(s,a,ds^+) = \mathbb E_\pi\left[\sum_{t=0}^{\infty}\gamma^t\delta_{s_{t+1}}(ds^+) \,\middle|\, s_0=s,\ a_0=a\right].
    \label{eq:M-definition}
\end{equation}
It has total mass $1/(1-\gamma)$ and satisfies
\begin{equation}
    M^\pi(s,a,ds^+) = P(ds^+\mid s,a) + \gamma\mathbb E_{\substack{s'\sim P(\cdot\mid s,a)\\a'\sim\pi(\cdot\mid s')}}\left[M^\pi(s',a',ds^+)\right].
    \label{eq:M-bellman}
\end{equation}

Assuming $M^\pi(s,a,\cdot)\ll\rho$, define
\[
    m^\pi(s,a,s^+) = \frac{dM^\pi(s,a,\cdot)}{d\rho}(s^+),
    \qquad
    q(s,a,s^+) = \frac{dP(\cdot\mid s,a)}{d\rho}(s^+).
\]
Equation~\eqref{eq:M-bellman} then becomes
\begin{equation}
    m^\pi(s,a,s^+) = q(s,a,s^+) + \gamma\mathbb E_{\substack{s'\sim P(\cdot\mid s,a)\\a'\sim\pi(\cdot\mid s')}}\left[m^\pi(s',a',s^+)\right].
    \label{eq:m-bellman}
\end{equation}

Standard FB models this unnormalized successor density with the bilinear factorization
\[
    m_z(s,a,s^+) \approx F(s,a,z)^\top B(s^+),
\]
and learns it through a squared Bellman objective. Density-FB retains the bilinear representation but changes the object being represented: rather than fitting the unnormalized successor measure, we model its normalized distribution over future states.


% ================================================================
\section{Density-FB}
\label{sec:density-fb}
% ================================================================

\subsection{Normalized successor density}

For a policy $\pi_z$, define
\begin{equation}
    \nu^{\pi_z} := (1-\gamma)M^{\pi_z}.
    \label{eq:nu-definition}
\end{equation}
Since $M^{\pi_z}$ has mass $1/(1-\gamma)$, $\nu^{\pi_z}(s,a,\cdot)$ is a probability distribution over future states. Its Bellman recursion is
\begin{equation}
    \nu^{\pi_z}(s,a,ds^+) = (1-\gamma)P(ds^+\mid s,a) + \gamma\mathbb E_{\substack{s'\sim P(\cdot\mid s,a)\\a'\sim\pi_z(\cdot\mid s')}}\left[\nu^{\pi_z}(s',a',ds^+)\right].
    \label{eq:nu-bellman}
\end{equation}

Assuming $\nu^{\pi_z}(s,a,\cdot)\ll\rho$, define its density ratio
\begin{equation}
    h_z^\pi(s,a,s^+) := \frac{d\nu^{\pi_z}(s,a,\cdot)}{d\rho}(s^+),
    \qquad
    \mathbb E_{s^+\sim\rho}\left[h_z^\pi(s,a,s^+)\right]=1.
    \label{eq:true-density-ratio}
\end{equation}

We estimate this density ratio through conditional density estimation. For a real-valued score $f(s,a,z,s^+)$, define
\begin{equation}
    h_f(s,a,z,s^+) =
    \frac{\exp f(s,a,z,s^+)}
    {\mathbb E_{\tilde s^+\sim\rho}\left[\exp f(s,a,z,\tilde s^+)\right]}.
    \label{eq:generic-density}
\end{equation}
For fixed $(s,a,z)$, consider
\begin{equation}
    \mathcal J_{\mathrm{ratio}}(f)
    =
    \mathbb E_{s^+\sim\nu^{\pi_z}(\cdot\mid s,a)}
    \left[f(s,a,z,s^+)\right]
    -
    \log\mathbb E_{s^+\sim\rho}
    \left[e^{f(s,a,z,s^+)}\right].
    \label{eq:ratio-objective}
\end{equation}
This objective is equal, up to an $f$-independent constant, to $-\operatorname{KL}\left( \nu^{\pi_z}(\cdot\mid s,a) \,\middle\|\, h_f(s,a,z,\cdot)\rho\right).$

Thus the normalized model recovers the successor density ratio at the population optimum. The corresponding score is the log density ratio up to an additive $(s,a,z)$-dependent constant, which is removed by normalization.


\subsection{Bilinear parameterization and learning objective}

We parameterize the log density ratio using the FB bilinear form
\begin{equation}
    f_\theta(s,a,z,s^+) =
    \frac{F_\theta(s,a,z)^\top B_\theta(s^+)}{\tau},
    \qquad
    h_\theta(s,a,z,s^+) =
    \frac{\exp f_\theta(s,a,z,s^+)}
    {\mathbb E_{\tilde s^+\sim\rho}
    \left[\exp f_\theta(s,a,z,\tilde s^+)\right]}.
    \label{eq:dfb-density}
\end{equation}
Unlike standard FB, the bilinear product therefore parameterizes a log-density score rather than the successor density directly.

Let $\bar F$ and $\bar B$ denote fixed target representations during a representation update, and let $\bar\pi=\sg(\pi)$. We define the target density ratio analogously as
\[
    \bar h(s',a',z,s^+)
    =
    \frac{
        \exp\left(\bar F(s',a',z)^\top\bar B(s^+)/\tau\right)
    }{
        \mathbb E_{\tilde s^+\sim\rho}
        \left[
            \exp\left(\bar F(s',a',z)^\top\bar B(\tilde s^+)/\tau\right)
        \right]
    }.
\]

Applying Equation~\eqref{eq:ratio-objective} to the normalized Bellman target in Equation~\eqref{eq:nu-bellman} gives the Density-FB population objective
\begin{equation}
\begin{aligned}
    \max_\theta\quad
    \mathcal J_{\mathrm{DFB}}
    ={}&
    (1-\gamma)
    \mathbb E_{\substack{
        z\sim p_Z\\
        (s,a,s')\sim\mathcal D
    }}
    \left[
        f_\theta(s,a,z,s')
    \right]
    \\
    &+
    \gamma
    \mathbb E_{\substack{
        z\sim p_Z\\
        (s,a,s')\sim\mathcal D\\
        a'\sim\bar\pi(\cdot\mid s',z)\\
        s^+\sim\rho
    }}
    \left[
        \sg\left[\bar h(s',a',z,s^+)\right]
        f_\theta(s,a,z,s^+)
    \right]
    \\
    &-
    \mathbb E_{\substack{
        z\sim p_Z\\
        (s,a)\sim\mathcal D_{SA}
    }}
    \left[
        \log
        \mathbb E_{s^+\sim\rho}
        \exp f_\theta(s,a,z,s^+)
    \right].
\end{aligned}
\label{eq:dfb-final}
\end{equation}

The first term anchors the density on the observed one-step successor. The second bootstraps the remaining discounted successor distribution from the target model. The final log-partition term normalizes the resulting score into a conditional density ratio.


\subsection{Finite-batch approximation}

For a minibatch $\{(s_i,a_i,s_i')\}_{i=1}^{n}$, we use the batch next states $\{s_j'\}_{j=1}^{n}$ as candidate future states. Define
\begin{equation}
    f_{ij}
    =
    \frac{F_\theta(s_i,a_i,z_i)^\top B_\theta(s_j')}{\tau},
    \qquad
    p_{ij}
    =
    \frac{\exp f_{ij}}
    {\sum_{k=1}^{n}\exp f_{ik}}.
    \label{eq:batch-density}
\end{equation}

For $a_i'\sim\pi(\cdot\mid s_i',z_i)$, treated as stop-gradient, the target network gives
\begin{equation}
    \bar p_{ij}'
    =
    \frac{
        \exp\left(
            \bar F(s_i',a_i',z_i)^\top\bar B(s_j')/\tau
        \right)
    }{
        \sum_{k=1}^{n}
        \exp\left(
            \bar F(s_i',a_i',z_i)^\top\bar B(s_k')/\tau
        \right)
    }.
    \label{eq:batch-target-density}
\end{equation}

The normalized Bellman recursion becomes the categorical target
\begin{equation}
    p_{ij}^{\mathrm{tgt}}
    =
    (1-\gamma)\mathbf 1[j=i]
    +
    \gamma\,\sg\left[\bar p_{ij}'\right],
    \qquad
    \mathcal L_{\mathrm{DFB}}
    =
    -\frac{1}{n}
    \sum_{i=1}^{n}
    \sum_{j=1}^{n}
    p_{ij}^{\mathrm{tgt}}\log p_{ij}.
    \label{eq:dfb-finite-pool}
\end{equation}

The batch softmax is therefore the finite-pool approximation of the normalized successor density in Equation~\eqref{eq:dfb-density}. For stability, we use two forward score heads with a shared backward representation and average their representation losses.


\subsection{Policy extraction}
\label{sec:policy-extraction}

The density parameterization changes how values are read from the representation. In standard FB, the linear successor factorization gives the reward embedding
\[
    z_r=\mathbb E_{s^+\sim\rho}\left[r(s^+)B(s^+)\right]
\]
and, under an exact factorization,
\begin{equation}
    \mathbb E_{\pi_{z_r}}
    \left[
        \sum_{t=0}^{\infty}\gamma^t r(s_{t+1})
        \,\middle|\,
        s_0=s,\ a_0=a
    \right]
    =
    F(s,a,z_r)^\top z_r.
    \label{eq:standard-fb-value}
\end{equation}

Density-FB no longer has this linear value readout because $F^\top B$ parameterizes a log density. Instead, we evaluate a reward directly under the predicted successor distribution. For any state reward $r$,
\begin{equation}
    q_r^{\pi_z}(s,a)
    :=
    \mathbb E_{s^+\sim\rho}
    \left[
        h_z^\pi(s,a,s^+)r(s^+)
    \right]
    =
    (1-\gamma)
    \mathbb E_{\pi_z}
    \left[
        \sum_{t=0}^{\infty}\gamma^t r(s_{t+1})
        \,\middle|\,
        s_0=s,\ a_0=a
    \right].
    \label{eq:true-normalized-value}
\end{equation}
Density-FB estimates this quantity as
\begin{equation}
    q_{\theta,r}(s,a;z)
    :=
    \mathbb E_{s^+\sim\rho}
    \left[
        h_\theta(s,a,z,s^+)r(s^+)
    \right]
    =
    \frac{
        \mathbb E_{s^+\sim\rho}
        \left[
            \exp f_\theta(s,a,z,s^+)r(s^+)
        \right]
    }{
        \mathbb E_{s^+\sim\rho}
        \left[
            \exp f_\theta(s,a,z,s^+)
        \right]
    }.
    \label{eq:model-reward-value}
\end{equation}

With a finite candidate pool $\{s_j^+\}_{j=1}^{M}$,
\begin{equation}
    q_{\theta,r}(s,a;z)
    \approx
    \sum_{j=1}^{M}
    p_j(s,a,z)r(s_j^+),
    \qquad
    p_j(s,a,z)
    =
    \frac{
        \exp f_\theta(s,a,z,s_j^+)
    }{
        \sum_k\exp f_\theta(s,a,z,s_k^+)
    }.
    \label{eq:finite-reward-value}
\end{equation}

\paragraph{Policy improvement.}
During training, the actor uses the reward $r_z$ associated with the sampled task index and optimizes $q_{\theta,z}:=q_{\theta,r_z}$. The representation parameters are frozen during the actor update, but gradients pass through the action input to $F$. Let
\begin{equation}
    s_q
    =
    \sg\left(
        \mathbb E
        \left|
            q_{\theta,z}
            -
            \mathbb E[q_{\theta,z}]
        \right|
        +
        \varepsilon
    \right).
    \label{eq:q-scale}
\end{equation}
We use the behavior-regularized actor objective
\begin{equation}
    \mathcal L_{\mathrm{actor}}
    =
    -
    \frac{
        \mathbb E_{s,z}
        \left[
            q_{\theta,z}(s,\pi(s,z))
        \right]
    }{
        s_q
    }
    +
    \beta\mathcal L_{\mathrm{BC}}(\pi).
    \label{eq:actor-centered}
\end{equation}
With two forward heads, the actor uses the minimum reward estimate across heads.

\paragraph{Goal-conditioned training.}
For the goal-conditioned experiments, $z$ is a goal $g$ sampled from $p_G$ and
\[
    r_z(\cdot)=r_g(\cdot)=\kappa(g,\cdot).
\]
The representation update conditions on $g$, while the corresponding reward is used only for policy improvement.

\paragraph{Downstream reward selection.}
For an unseen downstream reward $r$, we retain the learned policy family and evaluate candidate task indices using Equation~\eqref{eq:model-reward-value}. For a finite set $\mathcal Z_{\mathrm{cand}}$,
\begin{equation}
    J_{\theta,r}(z)
    :=
    \mathbb E_{s\sim\mu_{\mathrm{sel}}}
    \left[
        q_{\theta,r}(s,\pi(s,z);z)
    \right],
    \qquad
    \hat z_r
    \in
    \arg\max_{z\in\mathcal Z_{\mathrm{cand}}}
    J_{\theta,r}(z).
    \label{eq:policy-search}
\end{equation}
The selected policy $\pi_{\hat z_r}$ is then executed without updating $F$, $B$, or $\pi$. This selects the best estimated policy within the family covered during training; it does not imply optimality over policies outside that family.


% ================================================================
\section{Implementation}
\label{sec:implementation}
% ================================================================

Algorithm~\ref{alg:dfb} summarizes the complete training procedure.

\begin{algorithm}[h]
\caption{Density-FB}
\label{alg:dfb}
\begin{algorithmic}[1]
\Require dataset $\mathcal D$, task distribution $p_{\mathcal T}$, forward heads $F_\theta^{(1)},F_\theta^{(2)}$, shared backward representation $B_\theta$, targets $\bar F^{(1)},\bar F^{(2)},\bar B$, actor $\pi$, temperature $\tau$, coefficients $\beta,\varepsilon$, rates $\eta,\zeta$
\For{each training step}
    \State sample minibatch $\{(s_i,a_i,s_i')\}_{i=1}^{n}$
    \State sample $(z_i,r_{z_i})\sim p_{\mathcal T}$
    \For{$h\in\{1,2\}$}
        \State $p_{ij}^{(h)}
        \gets
        \mathrm{softmax}_j
        \left(
            F_\theta^{(h)}(s_i,a_i,z_i)^\top
            B_\theta(s_j')/\tau
        \right)$
    \EndFor
    \State $a_i'\gets\pi(s_i',z_i)$, stop-gradient
    \For{$h\in\{1,2\}$}
        \State $\bar p_{ij}'^{(h)}
        \gets
        \mathrm{softmax}_j
        \left(
            \bar F^{(h)}(s_i',a_i',z_i)^\top
            \bar B(s_j')/\tau
        \right)$
        \State $p_{ij}^{\mathrm{tgt},(h)}
        \gets
        \sg\left[
            (1-\gamma)\mathbf 1[j=i]
            +
            \gamma\bar p_{ij}'^{(h)}
        \right]$
    \EndFor
    \State update $F_\theta^{(1)},F_\theta^{(2)},B_\theta$ by minimizing
    \[
        -\frac{1}{2n}
        \sum_{h=1}^{2}
        \sum_{i,j}
        p_{ij}^{\mathrm{tgt},(h)}
        \log p_{ij}^{(h)}
    \]
    \State $a_i^\pi\gets\pi(s_i,z_i)$
    \For{$h\in\{1,2\}$}
        \State $p_{ij}^{\pi,(h)}
        \gets
        \mathrm{softmax}_j
        \left(
            F_\theta^{(h)}(s_i,a_i^\pi,z_i)^\top
            B_\theta(s_j')/\tau
        \right)$ with $F,B$ frozen
        \State $q_i^{(h)}
        \gets
        \sum_j
        p_{ij}^{\pi,(h)}
        r_{z_i}(s_j')$
    \EndFor
    \State $q_i\gets\min_h q_i^{(h)}$
    \State update $\pi$ using Equation~\eqref{eq:actor-centered}
    \State update target networks with rate $\zeta$
\EndFor
\Statex
\Statex \textbf{Deployment} (frozen $F,B,\pi$)
\State goal-conditioned task $g$: execute $\pi(s,g)$
\State downstream reward $r$: select $\hat z_r$ using Equation~\eqref{eq:policy-search}, then execute $\pi_{\hat z_r}$
\end{algorithmic}
\end{algorithm}


% ================================================================
\section{Evaluation}
\label{sec:evaluation}
% ================================================================

We evaluate three questions: whether the learned representation supports goal-conditioned control, whether the same frozen representation can evaluate rewards not used during training, and how the resulting zero-shot performance compares with published unsupervised RL methods. The successor representation is trained without downstream reward labels, while the policy family is co-trained using the task rewards sampled from $p_{\mathcal T}$. Unless stated otherwise, every model is trained for $500$k gradient steps, and intervals are $95\%$ $t$-based confidence intervals over three seeds.


\subsection{Goal-conditioned control}

We evaluate five goals per environment with $50$ episodes per goal. BC uses the same actor architecture with the value term removed, and GCIQL is our matched goal-conditioned control baseline.

\begin{table}[h]
    \centering
    \setlength{\tabcolsep}{5pt}
    \begin{tabular}{l cccc}
    \toprule
    Dataset & BC & FB & GCIQL & Density-FB \\
    \midrule
    \multicolumn{5}{l}{\emph{OGBench}}\\
    \texttt{cube-single} & $0.086\pm0.011$ & $0.420$ & $0.732\pm0.018$ & $\mathbf{0.915\pm0.098}$ \\
    \texttt{cube-double} & $0.023\pm0.014$ & $0.005\pm0.016$ & $0.377\pm0.111$ & $\mathbf{0.674\pm0.051}$ \\
    \texttt{scene} & $0.048\pm0.019$ & $0.076\pm0.057$ & $0.589\pm0.020$ & $\mathbf{0.641\pm0.123}$ \\
    \texttt{antmaze-medium} & $0.160\pm0.024$ & $0.593\pm0.250$ & $0.253\pm0.044$ & $\mathbf{0.676\pm0.188}$ \\
    \midrule
    \multicolumn{5}{l}{\emph{Metaworld} \texttt{play-resample}}\\
    \texttt{mw-pick-place} & $0.440\pm0.081$ & $0.627\pm0.052$ & $\mathbf{0.993\pm0.014}$ & $0.760\pm0.201$ \\
    \texttt{mw-push} & $0.540\pm0.000$ & $0.627\pm0.079$ & $0.733\pm0.198$ & $\mathbf{0.820\pm0.045}$ \\
    \midrule
    \multicolumn{5}{l}{\emph{Metaworld} \texttt{play-noise-hard-1m}}\\
    \texttt{mw-pick-place} & $0.340\pm0.023$ & $0.427\pm0.114$ & $\mathbf{0.520\pm0.060}$ & $0.460\pm0.045$ \\
    \texttt{mw-push} & $0.327\pm0.079$ & $0.340\pm0.060$ & $\mathbf{0.527\pm0.057}$ & $0.500\pm0.099$ \\
    \midrule
    \multicolumn{5}{l}{\emph{Metaworld} \texttt{play-resample-hard}}\\
    \texttt{mw-pick-place} & $0.367\pm0.124$ & $0.553\pm0.035$ & $\mathbf{0.967\pm0.035}$ & $0.573\pm0.372$ \\
    \texttt{mw-push} & $0.613\pm0.035$ & $0.607\pm0.111$ & $0.813\pm0.072$ & $\mathbf{0.827\pm0.094}$ \\
    \bottomrule
    \end{tabular}
    \caption{Native goal-conditioned success. All methods see the same datasets and evaluation goals.}
    \label{tab:main}
\end{table}


\subsection{Zero-shot reward transfer}

For zero-shot transfer, both the successor representation and policy family are frozen. The downstream task is specified by a reward $r$ that was not sampled from the training-task distribution. Density-FB evaluates each candidate policy through $q_{\theta,r}$ and selects its task index using Equation~\eqref{eq:policy-search}. We construct $\mathcal Z_{\mathrm{cand}}$ from $16{,}384$ dataset configurations and execute the selected policy without additional learning.

We consider three reward families defined from object coordinates and not used to train the representation: \emph{sparse}, the benchmark tolerance indicator around the five task goals; \emph{dense}, the same centers with distance shaping; and \emph{held-out regions}, the sparse reward centered at reachable dataset configurations that are not benchmark goals. Comparing sparse and dense rewards changes reward shaping while holding the target location fixed; comparing sparse and held-out regions changes the target location while holding the reward form fixed.

\begin{table}[h]
\centering
\small
\setlength{\tabcolsep}{5pt}
\begin{tabular}{llcc}
\toprule
Environment & Reward family & Plain FB & \DensityFB{} \\
\midrule
\texttt{cube-single} & sparse & $0.413 \pm 0.308$ & $\mathbf{0.888 \pm 0.091}$ \\
& dense & $0.488 \pm 0.397$ & $\mathbf{0.912 \pm 0.087}$ \\
& held-out regions & $0.090 \pm 0.114$ & $\mathbf{0.647 \pm 0.301}$ \\
\midrule
\texttt{cube-double} & sparse & $0.013 \pm 0.023$ & $\mathbf{0.573 \pm 0.126}$ \\
& dense & $0.037 \pm 0.011$ & $\mathbf{0.520 \pm 0.121}$ \\
& held-out regions & $0.000 \pm 0.000$ & $\mathbf{0.167 \pm 0.014}$ \\
\midrule
\texttt{scene} & sparse & $0.472 \pm 0.124$ & $\mathbf{0.640 \pm 0.111}$ \\
& dense & $0.413 \pm 0.030$ & $\mathbf{0.627 \pm 0.094}$ \\
& held-out regions & $0.053 \pm 0.080$ & $\mathbf{0.347 \pm 0.038}$ \\
\bottomrule
\end{tabular}
\caption{Zero-shot reward transfer (mean $\pm$ $95\%$ CI, three seeds). Plain FB uses its standard reward interface $z_r\propto\mathbb E_\rho[r(x)B(x)]$. Density-FB evaluates each covered candidate policy under the new reward and selects the highest-valued candidate according to Equation~\eqref{eq:policy-search}. Reward families, candidate pool, episode count, and success criterion are shared.}
\label{tab:zeroshot}
\end{table}


\subsection{Comparison with published zero-shot results}

Table~\ref{tab:published} places the sparse-reward results next to published state-based OGBench numbers. The comparison is not controlled: the published results use $1$M gradient steps, $10$ seeds, and $10$ episodes per task, whereas Density-FB uses $500$k steps, three seeds, and $25$ episodes. The published methods infer a task index by linear regression, whereas Density-FB selects among the learned goal-conditioned policy family using Equation~\eqref{eq:policy-search}.

\begin{table}[h]
\centering
\small
\begin{tabular}{lccc}
\toprule
Method & \texttt{cube-single} & \texttt{cube-double} & \texttt{scene} \\
\midrule
Laplacian (Wu et al., 2019) & $15.1 \pm 1.5$ & $2.0 \pm 0.4$ & $7.8 \pm 1.3$ \\
RLDP (Jajoo et al., 2026) & $19.8 \pm 2.4$ & $3.8 \pm 0.8$ & $11.6 \pm 1.6$ \\
ICVF (Ghosh et al., 2023) & $20.4 \pm 1.9$ & $5.0 \pm 0.8$ & $45.4 \pm 2.3$ \\
BYOL & $22.0 \pm 3.2$ & $4.4 \pm 0.7$ & $15.4 \pm 1.4$ \\
TD-JEPA (Bagatella et al., 2026) & $34.2 \pm 2.9$ & $3.6 \pm 0.8$ & $38.4 \pm 1.4$ \\
FB (Touati \& Ollivier, 2021) & $49.6 \pm 3.8$ & $2.6 \pm 0.4$ & $12.8 \pm 1.6$ \\
HILP (Park et al., 2024) & $74.2 \pm 3.5$ & $20.0 \pm 2.7$ & $43.8 \pm 1.9$ \\
BYOL-$\gamma$ & $79.4 \pm 2.8$ & $2.6 \pm 0.7$ & $14.4 \pm 2.3$ \\
\midrule
\DensityFB{} & $\mathbf{88.8 \pm 9.1}$ & $\mathbf{57.3 \pm 12.6}$ & $\mathbf{64.0 \pm 11.1}$ \\
\bottomrule
\end{tabular}
\caption{Zero-shot success rate (\%) on state-based OGBench. Published rows use $1$M gradient steps, $10$ seeds, and $10$ episodes per task. Density-FB uses $500$k steps, three seeds, and $25$ episodes per task. Density-FB confidence intervals are $t$-based over three seeds.}
\label{tab:published}
\end{table}


% ================================================================
\appendix
\section{Standard FB Bellman Objective}
\label{app:standard-fb}
% ================================================================

For completeness, standard FB can be obtained by fitting the unnormalized density in Equation~\eqref{eq:m-bellman} with a squared $L_2(\rho)$ Bellman residual. For a fixed target density $\bar m$,
\begin{equation}
    \mathcal L_M
    =
    \frac12
    \mathbb E_{\substack{
        (s,a,s')\sim\mathcal D\\
        a'\sim\pi(\cdot\mid s')\\
        s^+\sim\rho
    }}
    \left[
        m(s,a,s^+)
        -
        q(s,a,s^+)
        -
        \gamma\bar m(s',a',s^+)
    \right]^2.
    \label{eq:standard-fb-square}
\end{equation}

Using
\[
    \mathbb E_{s^+\sim\rho}
    \left[
        m(s,a,s^+)q(s,a,s^+)
    \right]
    =
    \mathbb E_{s'\sim P(\cdot\mid s,a)}
    \left[
        m(s,a,s')
    \right],
\]
and dropping terms independent of $m$ gives
\begin{equation}
\begin{aligned}
    \max_m\quad
    \mathcal J_M
    ={}&
    \mathbb E_{(s,a,s')\sim\mathcal D}
    \left[
        m(s,a,s')
    \right]
    \\
    &+
    \gamma
    \mathbb E_{\substack{
        (s,a,s')\sim\mathcal D\\
        a'\sim\pi(\cdot\mid s')\\
        s^+\sim\rho
    }}
    \left[
        m(s,a,s^+)\bar m(s',a',s^+)
    \right]
    \\
    &-
    \frac12
    \mathbb E_{\substack{
        (s,a,s')\sim\mathcal D\\
        s^+\sim\rho
    }}
    \left[
        m(s,a,s^+)^2
    \right].
\end{aligned}
\label{eq:standard-fb-objective}
\end{equation}

Substituting $m_z(s,a,s^+)=F(s,a,z)^\top B(s^+)$ yields the usual FB representation objective. The difference introduced by Density-FB is therefore not the bilinear score itself, but the decision to model the normalized successor density through an exponential normalization rather than regress the unnormalized successor measure directly.


% ================================================================
\section{Optional Reward-Dependent Policy Index}
\label{app:reward-index}
% ================================================================

Density normalization removes the standard FB identity
\[
    \int
    r(s^+)
    F(s,a,z)^\top B(s^+)
    \,d\rho(s^+)
    =
    F(s,a,z)^\top
    \mathbb E_\rho[rB],
\]
because $F^\top B$ now appears inside a normalized exponential. In particular, the density is invariant to additive shifts $B(s^+)\mapsto B(s^+)+b$, whereas the raw moment $\mathbb E_\rho[rB]$ is not.

If a reward-dependent index is required, one possible centered construction uses reward evaluations $\{(s_j^+,r_j)\}_{j=1}^{M}$ with $s_j^+\sim\rho$. Let
\[
    \bar r
    =
    \frac1M\sum_j r_j,
    \qquad
    \mu_B
    =
    \mathbb E_\rho[B],
    \qquad
    C_B
    =
    \mathbb E_\rho
    \left[
        (B-\mu_B)(B-\mu_B)^\top
    \right].
\]
Then, for $\lambda>0$,
\begin{equation}
    z_r
    =
    \sqrt d\,
    \normalize
    \left[
        (C_B+\lambda I)^{-1}
        \frac1M
        \sum_{j=1}^{M}
        (B(s_j^+)-\mu_B)(r_j-\bar r)
    \right].
    \label{eq:centered-dual}
\end{equation}

This estimator is invariant to additive shifts of $B$ and reward shifts, and after normalization to positive reward scaling. It requires an actor trained on the corresponding latent index space and is not used for the goal-search results reported in Section~\ref{sec:evaluation}.


\end{document}